\documentclass[a4paper,12pt]{article}

\usepackage[french]{babel}
\usepackage[T1]{fontenc}
\usepackage[utf8]{inputenc}
\usepackage{amsmath, amssymb}
\usepackage{geometry}
\usepackage{graphicx}
\usepackage{array}
\usepackage{enumitem}
\usepackage{tkz-tab}
\usepackage{tikz}
\usepackage{pgfplots}
\pgfplotsset{compat=1.18}
\usepgfplotslibrary{statistics}

\geometry{top=2cm, bottom=2cm, left=2cm, right=2cm}

\begin{document}

\begin{center}\Large\textbf{Correction - Épreuve anticipée de maths - Sujet 2}\end{center}

\section*{Partie 1 : Automatismes (12 questions -- 6 points)}

\begin{enumerate}

\item \textbf{Réponse :} \boxed{A} \(-2\)

\[
M = -1 + \frac{5 - 1}{2 \times (-2)} = -1 + \frac{4}{-4} = -1 - 1 = -2
\]

\item \textbf{Réponse :} \boxed{D} uniquement \(k_2\) et \(k_3\)

\[
k_1(x) = (2x-1)^2 - (x+1)(2x-1) = (2x-1)[(2x-1)-(x+1)] = (2x-1)(x-2)
\]
\[
k_1(x) = 2x^2 -5x +2 \text{ (non affine)}
\]
\[
k_2(x) = 3x - 5 + \frac{x}{2} = \frac{6x + x}{2} - 5 = \frac{7x}{2} - 5 \text{ (affine)}
\]
\[
k_3(x) = \frac{4x^2 - (4x^2 - 12x + 9)}{2} = \frac{4x^2 - 4x^2 + 12x - 9}{2} = \frac{12x - 9}{2} = 6x - 4,5 \text{ (affine)}
\]

\item \textbf{Réponse :} \boxed{C} \(10\)
\item \textbf{Réponse :} \boxed{B} \(33\%\)

\[
\frac{85}{260} \times 100 \approx 32,69\% \approx 33\%
\]

\item \textbf{Réponse :} \boxed{B} Dans la série 2, au moins \(75\%\) des valeurs sont inférieures ou égales à 13.

Analyse :
- A : Médiane série 1 = 9 (vrai, mais pas la seule affirmation vraie)
- B : Q3 = 13, donc 75\% des valeurs \(\leq\) 13 (vrai)
- C : Étendue série 1 = 10, série 2 = 16 (faux)
- D : Écart interquartile série 1 = 6, série 2 = 6 (faux, ils sont égaux)

\item \textbf{Réponse :} \boxed{B} \(10\,000\)

\[
(3\pi)^2 \times 120 = 9\pi^2 \times 120 = 1080 \times \pi^2
\]
\(\pi^2 \approx 9,86\), donc \(1080 \times 9,86 \approx 10\,648 \approx 10^4\)

\item \textbf{Réponse :} \boxed{C} \(16\)

\[
\frac{4^3 \times 3^2}{6^2} = \frac{64 \times 9}{36} = \frac{576}{36} = 16
\]

\item \textbf{Réponse :} \boxed{A}

\(f(x) = 3x - 1\) s'annule en \(x = \frac{1}{3}\), coefficient directeur \(3 > 0\) donc négative avant, positive après.

\item \textbf{Réponse :} \boxed{A} \(0,6\) g

Diminuer de 25\% revient à multiplier par \(0,75\) : \(0,8 \times 0,75 = 0,6\)

\item \textbf{Réponse :} \boxed{B} \(d = \sqrt{\dfrac{G m_1 m_2}{F}}\)

\[
F = G\frac{m_1 m_2}{d^2} \implies d^2 = \frac{G m_1 m_2}{F} \implies d = \sqrt{\frac{G m_1 m_2}{F}}
\]

\item \textbf{Réponse :} \boxed{D} \(25x^2 + 40x + 16\)

\[
(5x+4)^2 = 25x^2 + 2 \times 5x \times 4 + 16 = 25x^2 + 40x + 16
\]

\item \textbf{Réponse :} \boxed{B} \(0,6\)

\(P(\text{Allemand} \mid \text{Fille}) = \dfrac{30}{50} = 0,6 = \dfrac{3}{5}\)

\end{enumerate}

\section*{Partie 2 : Exercices (14 points)}

\subsection*{Exercice 1 (7 points)}

\begin{enumerate}
\item \textbf{Volume au printemps 2027 :}
\[
u_1 = 0,8 \times u_0 + 40 = 0,8 \times 1000 + 40 = 800 + 40 = 840 \text{ m}^3
\]

\item \textbf{Justification de la relation :}
\begin{itemize}
\item Perte de 20\% : on multiplie par \(0,8\)
\item On ajoute ensuite 40 m³
\item Donc \(u_{n+1} = 0,8 u_n + 40\)
\end{itemize}

\item \textbf{La suite n'est pas arithmétique :}
\[
u_1 - u_0 = 840 - 1000 = -160
\]
\[
u_2 = 0,8 \times 840 + 40 = 672 + 40 = 712
\]
\[
u_2 - u_1 = 712 - 840 = -128 \neq -160
\]
La différence n'est pas constante, donc \((u_n)\) n'est pas arithmétique.

\item \textbf{Étude de la suite auxiliaire :}
\begin{enumerate}
\item \(v_n = u_n - 200\)
\[
v_{n+1} = u_{n+1} - 200 = (0,8 u_n + 40) - 200 = 0,8 u_n - 160
\]
\[
v_{n+1} = 0,8(u_n - 200) = 0,8 v_n
\]
Donc \((v_n)\) est géométrique de raison \(q = 0,8\).

\item \(v_0 = u_0 - 200 = 1000 - 200 = 800\)
\[
v_n = v_0 \times q^n = 800 \times (0,8)^n
\]

\item \(u_n = v_n + 200 = 800 \times 0,8^n + 200\)
\end{enumerate}

\item \textbf{Analyse des valeurs :}
\begin{enumerate}
\item D'après le tableau, lorsque \(n\) augmente, \(u_n\) diminue et se rapproche de 200. On peut conjecturer que la suite \((u_n)\) converge vers 200 lorsque \(n\) tend vers \(+\infty\).
\item Le volume d'eau dans le bassin se stabilise autour de 200 m³ au bout de plusieurs années. Le bassin atteint un équilibre où les pertes sont compensées par l'apport annuel.
\end{enumerate}

\end{enumerate}

\subsection*{Exercice 2 (7 points)}

\begin{enumerate}
\item \textbf{Calcul de la dérivée :}
\[
f(x) = \frac{3x - 3}{2x + 1}
\]
Forme \(\frac{u}{v}\) avec \(u(x) = 3x - 3\), \(u'(x) = 3\), \(v(x) = 2x + 1\), \(v'(x) = 2\).
\[
f'(x) = \frac{u'v - uv'}{v^2} = \frac{3(2x+1) - (3x-3) \times 2}{(2x+1)^2}
\]
\[
f'(x) = \frac{6x + 3 - (6x - 6)}{(2x+1)^2} = \frac{6x + 3 - 6x + 6}{(2x+1)^2} = \frac{9}{(2x+1)^2}
\]

\item \textbf{Variations de \(f\) :}
\[
f'(x) = \frac{9}{(2x+1)^2} > 0 \text{ pour tout } x \neq -\frac{1}{2}
\]
Donc \(f\) est strictement croissante sur \(]-\infty ; -\frac{1}{2}[\) et sur \(]-\frac{1}{2} ; +\infty[\).

\item \textbf{Tangentes de coefficient directeur 1 :}
On résout \(f'(x) = 1\) :
\[
\frac{9}{(2x+1)^2} = 1 \iff (2x+1)^2 = 9
\]
\[
2x+1 = 3 \quad \text{ou} \quad 2x+1 = -3
\]
\[
2x = 2 \implies x = 1 \quad \text{ou} \quad 2x = -4 \implies x = -2
\]
Deux solutions distinctes, donc exactement deux tangentes.

\item \textbf{Équations des tangentes :}
Formule : \(y = f'(x_0)(x - x_0) + f(x_0)\)

\noindent Pour \(x_0 = 1\) :
\[
f(1) = \frac{3 \times 1 - 3}{2 \times 1 + 1} = \frac{0}{3} = 0
\]
\[
y = 1 \times (x - 1) + 0 = x - 1
\]

\noindent Pour \(x_0 = -2\) :
\[
f(-2) = \frac{3 \times (-2) - 3}{2 \times (-2) + 1} = \frac{-6 - 3}{-4 + 1} = \frac{-9}{-3} = 3
\]
\[
y = 1 \times (x + 2) + 3 = x + 5
\]

Les deux tangentes ont pour équations : \(y = x - 1\) et \(y = x + 5\).

\end{enumerate}

\vspace{1cm}
\textbf{Barème indicatif :}
\begin{itemize}
\item Automatismes : 6 points (0,5 pt par question)
\item Exercice 1 : 7 points
  \begin{itemize}
  \item 1) 1 pt
  \item 2) 1 pt
  \item 3) 1 pt
  \item 4a) 1 pt
  \item 4b) 1 pt
  \item 4c) 1 pt
  \item 5a) 0,5 pt
  \item 5b) 0,5 pt
  \end{itemize}
\item Exercice 2 : 7 points
  \begin{itemize}
  \item 1) 2 pts
  \item 2) 1,5 pts
  \item 3) 1,5 pts
  \item 4) 2 pts
  \end{itemize}
\end{itemize}

\textbf{Note concernant la question 3 des automatismes :} 
Les données du tableau conduisent à \(x = 6,2\) qui n'apparaît pas dans les propositions. Cette question sera neutralisée et tous les élèves obtiendront les 0,5 point.

\end{document}